\subsection{中心化子}

回忆，群的两个元素 \(x, y\) 称为可交换的，是指 \((x, y) = e\)、\((\operatorname{Int} x)y = y\) 或 \((\operatorname{Int} y)x = x\)；李代数的两个元素 \(a, b\) 称为可交换的，是指 \([a, b] = 0\)、\((\operatorname{ad} a).b = 0\) 或 \((\operatorname{ad} b).a = 0\)。设 \(G\) 是李群，\(x \in G\)、\(a \in L(G)\)；若 \(x\)\(a\)\((\operatorname{Ad} x).a = a\)，即在 \(xa = ax\) 中 \(T(G)\)，则称 x 与 a 可交换。

设 \(G\)\(g\)\(G\)\(a\)\(g\)\(Z_{G}(A)\) 是李群，g 是其李代数，A 是 G 的子集，a 是 g 的子集。令 \(Z_{G}(a)\)（分别地，\(G\)\(a\)\(G\)）表示与 A（分别地，a）中的所有元素可交换的 G 中元素的集合。它是 G 的闭子群。令 \(\mathfrak{z}_{\mathfrak{g}}(A)\)（分别地，\(\mathfrak{z}_{\mathfrak{g}}(a)\)\(g\)\(a\)\(g\)）表示与 A（分别地，a）中的所有元素可交换的 g 中元素的集合。它是 g 的闭李子代数。

命题 7. 设 \(G\)\(g\)\(a\)\(g\)\(Z_G(a)\) 是有限维李群，g 是其李代数，a 是 g 的子集。则 \(G\) 是 G 的李子群，其李代数为 \(\mathfrak{z}_{\mathfrak{g}}(a)\)。

这由 § 3，命题 44 以及命题 39 的推论 2 得到。

命题 8. 设 \(\mathbf{G}\) 是有限维实或复李群，\(\mathfrak{g}\) 是其李代数，\(\mathbf{A}\) 是 \(\mathbf{G}\) 的子集。则 \(Z_{\mathbf{G}}(\mathbf{A})\) 是 \(\mathbf{G}\) 的李子群，其李代数为 \(\mathfrak{z}_{\mathfrak{g}}(\mathbf{A})\)。

假设 A 只含有一个点 \(a\)。则 \(Z_{\mathbf{G}}(\mathbf{A})\) 是 Int \(a\) 的不动点集合；因此 \(Z_{\mathbf{G}}(\mathbf{A})\) 是 \(\mathbf{G}\) 的李子群，而 \(\mathbf{L}(Z_{\mathbf{G}}(\mathbf{A}))\) 是 Ad \(a\) 的不动点集合，即 \(\mathfrak{z}_{\mathfrak{g}}(\mathbf{A})\)（§ 3，第 8 号，命题 29 的推论 1）。一般情形由 § 6，第 2 号，命题 1 的推论 3 得到。

命题 9. 设 G 是有限维实或复李群，g 是其李代数，A 是 G 的积分子群，且 \(\mathfrak{a} = \mathrm{L}(\mathbf{A})\)。则 \(\mathrm{Z}_{\mathbf{G}}(\mathbf{A}) = \mathrm{Z}_{\mathbf{G}}(\mathfrak{a})\)、\(\mathfrak{z}_{\mathfrak{g}}(\mathbf{A}) = \mathfrak{z}_{\mathfrak{g}}(\mathfrak{a})\)，并且 \(\mathrm{Z}_{\mathbf{G}}(\mathbf{A})\) 是 G 的李子群，其李代数为 \(\mathfrak{z}_{\mathfrak{g}}(\mathfrak{a})\)。令 \(x \in G\)。则

\[
\begin{array}{l l} x \in Z _ {\mathrm{G}} (\mathrm{A}) \Leftrightarrow \mathrm{A} \subset Z _ {\mathrm{G}} (\{x \}) \\ \quad \Leftrightarrow \mathfrak {a} \subset \mathrm{L} (Z _ {\mathrm{G}} (\{x \})) & (\S 6, \text { Corollary   2   to   Proposition   3 }) \\ \quad \Leftrightarrow \mathfrak {a} \subset \mathfrak {z}_ {\mathfrak {g}} (\{x \}) & (\text { Proposition   8 }) \\ \quad \Leftrightarrow x \in Z _ {\mathrm{G}} (\mathfrak {a}) \end{array}
\]

从而 \(Z_{\mathbf{G}}(\mathbf{A}) = Z_{\mathbf{G}}(\mathfrak{a})\)。令 \(u \in \mathfrak{g}\)。则

\[
\begin{array}{r l} u \in \mathfrak{z}_{\mathfrak{g}}(\mathfrak{a}) & \Leftrightarrow \mathfrak{a} \subset Z_{\mathrm{G}} (\{u\}) \\ & \Leftrightarrow \mathfrak{a} \subset \mathrm{L} (Z_{\mathrm{G}} (\{u\})) \quad (\S 6, \text { Corollary   2   to   Proposition   3 }) \\ & \Leftrightarrow \mathfrak{a} \subset \mathfrak{z}_{\mathfrak{g}} (\{u\}) \quad (\text { Proposition   7 }) \\ & \Leftrightarrow u \in \mathfrak{z}_{\mathfrak{g}} (\mathfrak{a}) \end{array}
\]

从而 \(\mathfrak{z}_{\mathfrak{g}}(\mathbf{A}) = \mathfrak{z}_{\mathfrak{g}}(\mathfrak{a})\)。最后一个断言由命题 7 或命题 8 得到。

